%! TEX root = ../master.tex
\begin{theorem}[Reflective property of the ellipse]\label{thm:ellipse-reflection}
  A ray leaving one focus of an ellipse is reflected by the ellipse through
  the other focus.
\end{theorem}

\begin{proof}
  Take the ellipse of \cref{thm:ellipse-standard} with $c > 0$ (if $c = 0$
  it is a circle, and every ray from the centre returns along itself). Let
  $P(x_1, y_1)$ be on the ellipse.

  If $y_1 = 0$, then $P = (\pm a, 0)$, the tangent there is vertical, and the
  ray runs along the $x$-axis, so it returns along the axis through both
  foci. Suppose now $y_1 \neq 0$ and, for the moment, $x_1 \neq \pm c$, so
  that neither focal line is vertical.

  \emph{Slopes.} Differentiating $x^2/a^2 + y^2/b^2 = 1$ implicitly gives
  $2x/a^2 + (2y/b^2)\,y' = 0$. So the tangent at $P$ has slope
  \[
    m = -\frac{b^2 x_1}{a^2 y_1},
  \]
  and the lines $F_1P$ and $F_2P$ have slopes
  \[
    m_1 = \frac{y_1}{x_1 + c}, \qquad m_2 = \frac{y_1}{x_1 - c}.
  \]
  Below we use $b^2 x_1^2 + a^2 y_1^2 = a^2 b^2$ (because $P$ is on the
  ellipse) and $a^2 - b^2 = c^2$.

  \emph{Incoming angle.} Let $\alpha$ be the angle from $F_1P$ to the tangent.
  By \cref{thm:angle-between-lines},
  $\tan\alpha = (m - m_1)/(1 + m m_1)$. Multiply the numerator and the
  denominator by $a^2 y_1 (x_1 + c)$:
  \[
    \begin{aligned}
      \text{numerator}
      &= -b^2 x_1 (x_1 + c) - a^2 y_1^2
       = -(b^2 x_1^2 + a^2 y_1^2) - b^2 c x_1
       = -b^2 (a^2 + c x_1), \\
      \text{denominator}
      &= a^2 y_1 (x_1 + c) - b^2 x_1 y_1
       = y_1 \bigl( (a^2 - b^2) x_1 + a^2 c \bigr)
       = c y_1 (a^2 + c x_1).
    \end{aligned}
  \]
  Since $\abs{x_1} \leq a$ and $c < a$, we have $a^2 + c x_1 > 0$. So the
  denominator is not zero (the lines are not perpendicular), and
  \[
    \tan\alpha = -\frac{b^2}{c y_1}.
  \]

  \emph{Outgoing angle.} Let $\beta$ be the angle from the tangent to $F_2P$,
  so $\tan\beta = (m_2 - m)/(1 + m m_2)$. Multiply the numerator and the
  denominator by $a^2 y_1 (x_1 - c)$:
  \[
    \begin{aligned}
      \text{numerator}
      &= a^2 y_1^2 + b^2 x_1 (x_1 - c)
       = b^2 (a^2 - c x_1), \\
      \text{denominator}
      &= a^2 y_1 (x_1 - c) - b^2 x_1 y_1
       = -c y_1 (a^2 - c x_1),
    \end{aligned}
  \]
  with $a^2 - c x_1 > 0$. Hence
  \[
    \tan\beta = -\frac{b^2}{c y_1} = \tan\alpha.
  \]

  Both angles lie in $[0, \pi)$, where $\tan$ is one-to-one, so
  $\alpha = \beta$. So $F_1P$ and $F_2P$ make equal angles with the tangent,
  one on each side: reflecting the line $F_1P$ in the tangent gives the line
  $F_2P$. This is the law of reflection (the angle of incidence equals the
  angle of reflection), so the reflected ray passes through $F_2$.

  When $x_1 = \pm c$, one focal line is vertical and the slope formula does
  not apply. The result still holds there by continuity: both angles depend
  continuously on $P$, and they agree at every nearby point of the ellipse.
\end{proof}
